\chapter{Problemas frecuentes}
\label{cap:problemas-frecuentes}

\begin{procedimiento}{Objetivo del capitulo}{Todos}
Resolver problemas operativos comunes con una accion inmediata y una referencia al capitulo donde se explica el flujo completo.
\end{procedimiento}

Use esta guia cuando algo no coincida con lo esperado. No sustituye la validacion de la organizacion ni el soporte tecnico. Si el problema involucra resultados oficiales, correos reales o cambios en Django Admin, detenga la operacion y verifique antes de repetir.

\section{Guia rapida}

\begin{longtable}{p{3.0cm}p{3.8cm}p{4.4cm}p{1.9cm}}
\toprule
Problema & Causa probable & Accion inmediata & Capitulo\\
\midrule
\endhead
La pagina no carga & Red, servidor no disponible, ruta incorrecta o sesion expirada. & Recargue una vez, pruebe otra pantalla y reporte si el fallo persiste. & \cref{cap:acceso-navegacion}\\
No puede iniciar sesion & Credenciales incorrectas, cuenta sin permiso o sesion no autorizada. & Verifique usuario autorizado. No comparta contrasena. & \cref{cap:acceso-navegacion}\\
El formulario no se envia & Campos invalidos, duplicados o edicion activa ausente. & Lea errores junto a los campos y corrija. & \cref{cap:registro-participantes}\\
Faltan campos obligatorios & Datos de institucion, responsable, robot o integrante incompletos. & Complete campos marcados y vuelva a enviar. & \cref{cap:registro-participantes}\\
El enlace de confirmacion no funciona & Token inexistente, enlace copiado incompleto o registro no disponible. & No confirme por otra via; solicite enlace valido. & \cref{cap:registro-participantes}\\
Un participante no aparece & Asistencia no confirmada, equipo no sincronizado, filtro equivocado o division incorrecta. & Revise asistencia, equipo y division antes de crear duplicados. & \cref{cap:gestion-organizacion}\\
La asistencia no se refleja & Confirmacion pendiente, solicitud no enviada o sincronizacion no ejecutada. & Revise registro y accion de sincronizacion autorizada. & \cref{cap:gestion-organizacion,cap:comunicaciones}\\
Los cambios no se guardaron & Seleccion visual sin guardar, error AJAX o cancelacion de confirmacion. & Recargue y verifique estado persistido antes de repetir. & \cref{cap:fase-grupos}\\
El boton no responde & No hay cambios, validacion pendiente, boton deshabilitado o fase incompleta. & Lea mensaje visible y revise precondiciones. & \cref{cap:estados-mensajes}\\
Error de red durante AJAX & Conexion interrumpida o respuesta inesperada. & No asuma guardado. Recargue y compruebe el resultado. & \cref{cap:estados-mensajes}\\
La seleccion desaparece tras recargar & No se habia guardado o la operacion fallo. & Repita seleccion solo despues de confirmar el estado actual. & \cref{cap:fase-grupos,cap:eliminatorias-repechaje}\\
No se habilita la fase siguiente & Faltan resultados guardados en grupo o batalla. & Complete y guarde todos los resultados de la fase actual. & \cref{cap:eliminatorias-repechaje}\\
Resultado visual inesperado & Estado recalculado, correccion previa o fase posterior afectada. & Compare con acta y revise historial o modal de participante. & \cref{cap:correcciones-manuales}\\
Aparece confirmacion de correccion & El cambio afecta progreso posterior. & Cancele si no hay autorizacion; si acepta, revise fases posteriores. & \cref{cap:correcciones-manuales}\\
Repechaje sin candidatos & No hay suficientes participantes elegibles o la fase no esta cerrada. & Cierre resultados actuales y revise recomendacion. & \cref{cap:eliminatorias-repechaje}\\
El podio no se actualiza & Seleccion incompleta, posiciones duplicadas o error al guardar. & Revise campeon, segundo y tercero; luego guarde y verifique. & \cref{cap:fases-finales-podio}\\
Previsualizacion no coincide & Marcadores DOCX faltantes o datos extra distintos. & Revise marcadores detectados y plantilla activa. & \cref{cap:comunicaciones}\\
Lote en error & SMTP bloqueado, destinatario invalido, plantilla no valida o envio real sin confirmacion. & Revise historial y corrija antes de repetir. & \cref{cap:comunicaciones}\\
Admin muestra datos inesperados & Filtro activo, edicion equivocada, dato modificado directamente o busqueda incompleta. & Quite filtros, revise edicion y no cree duplicados. & \cref{cap:django-admin}\\
Interfaz parece desactualizada & Cache del navegador o pagina abierta antes de un cambio. & Recargue la pagina y verifique que no haya cambios sin guardar. & \cref{cap:acceso-navegacion}\\
\bottomrule
\end{longtable}

\section{Diferenciar responsabilidades}

\begin{tabularx}{\linewidth}{>{\bfseries}p{3.2cm}X}
\toprule
Tipo de respuesta & Alcance\\
\midrule
Usuario publico & Puede corregir campos de registro, revisar reglas, volver a abrir enlaces y solicitar orientacion si el enlace no funciona. No debe acceder al panel interno.\\
Operador & Puede guardar resultados, revisar fases, simular comunicaciones, consultar historial y detener la operacion si hay duda.\\
Superusuario Admin & Puede revisar modelos y acciones administrativas autorizadas; no debe usar Admin para reemplazar flujos normales si existe una pantalla operativa.\\
Soporte tecnico & Debe intervenir cuando hay error repetido, datos incoherentes, fallo de servidor, problema de permisos o necesidad de respaldo segun politica aprobada.\\
\bottomrule
\end{tabularx}

\section{Antes de solicitar soporte}

Prepare una descripcion breve con:
\begin{itemize}
    \item Pantalla o ruta visible.
    \item Fecha y hora aproximada.
    \item Division, fase, grupo, batalla o lote afectado.
    \item Accion realizada.
    \item Mensaje mostrado.
    \item Captura sin correos, documentos, tokens ni credenciales.
    \item Resultado esperado.
    \item Resultado obtenido.
\end{itemize}

No envie contrasena, token de asistencia, claves secretas, credenciales SMTP ni datos personales que no sean necesarios para diagnosticar el caso. Contacto o canal de soporte: \campointerfaz{PENDIENTE DE DEFINICION INSTITUCIONAL}.

\section{Reglas de seguridad ante incidentes}

\begin{itemize}
    \item No repita envios reales para probar si el problema se arreglo.
    \item No acepte confirmaciones sensibles si no entiende el impacto.
    \item No cree registros duplicados para que un participante aparezca.
    \item No use Django Admin para corregir resultados si existe un flujo operativo documentado.
    \item No continue una fase si el resultado visible no coincide con el acta o decision oficial.
\end{itemize}
